\section{Реализация программы}

Реализация программы представлена в листинге \ref{lst:prog}.

\begin{lstlisting}[caption={Реализация программы <<Предки>>}, label={lst:prog}]
find_grandparent(GP, GPGender, PGender, Child) :-
  parent(GP, Parent),
  gender(GP, GPGender),
  parent(Parent, Child),
  gender(Parent, PGender).

gender(male, petr).
gender(male, fedor).
gender(male, alexey).
gender(male, ivan).

gender(female, anna).
gender(female, olga).
gender(female, maria).

parent(petr, maria).
parent(anna, maria).
parent(fedor, alexey).
parent(olga, alexey).
parent(maria, ivan).
parent(alexey, ivan).

max2_no_cut(X, Y, X) :- X >= Y.
max2_no_cut(X, Y, Y) :- X < Y.
max2_cut(X, Y, X) :- X >= Y, !.
max2_cut(_, Y, Y).

max3_no_cut(X, Y, Z, X) :- X >= Y, X >= Z.
max3_no_cut(X, Y, Z, Y) :- Y > X, Y >= Z.
max3_no_cut(X, Y, Z, Z) :- Z > X, Z > Y.
max3_cut(X, Y, Z, X) :- X >= Y, X >= Z, !.
max3_cut(_, Y, Z, Y) :- Y >= Z, !.
max3_cut(_, _, Z, Z).
\end{lstlisting}

Форма вопроса вида \ref{lst:query1} по имени субъекта определить всех его бабушек.
\begin{lstlisting}[caption={Форма задания вопроса 1}, label={lst:query1}]
?- find_grandparent(GP, female, _, name).
\end{lstlisting}

Форма вопроса вида \ref{lst:query2} по имени субъекта определить всех его дедушек.
\begin{lstlisting}[caption={Форма задания вопроса 2}, label={lst:query2}]
?- find_grandparent(GP, male, _, name).
\end{lstlisting}

Форма вопроса вида \ref{lst:query3} по имени субъекта определить всех его бабушек и дедушек.
\begin{lstlisting}[caption={Форма задания вопроса 3}, label={lst:query3}]
?- find_grandparent(GP, _, _, name).
\end{lstlisting}

Форма вопроса вида \ref{lst:query4} по имени субъекта определить его бабушку по материнской линии.
\begin{lstlisting}[caption={Форма задания вопроса 4}, label={lst:query4}]
?- find_grandparent(GP, female, female, name).
\end{lstlisting}

Форма вопроса вида \ref{lst:query5} по имени субъекта определить его бабушку и дедушку по материнской линии.
\begin{lstlisting}[caption={Форма задания вопроса 5}, label={lst:query5}]
?- find_grandparent(GP, _, female, name).
\end{lstlisting}

\section{Варианты вопросов и ответов на них}

\begin{lstlisting}[caption={Вопрос и ответ №1}, label={lst:q1}]
?- find_grandparent(GP, female, _, ivan).
GP = anna ;
GP = olga .
\end{lstlisting}

\begin{lstlisting}[caption={Вопрос и ответ №2}, label={lst:q2}]
?- find_grandparent(GP, male, _, ivan).
GP = petr ;
GP = fedor .
\end{lstlisting}

\begin{lstlisting}[caption={Вопрос и ответ №3}, label={lst:q3}]
?- find_grandparent(GP, _, _, ivan).
GP = petr ;
GP = anna ;
GP = fedor ;
GP = olga .
\end{lstlisting}

\begin{lstlisting}[caption={Вопрос и ответ №4}, label={lst:q4}]
?- find_grandparent(GP, female, female, ivan).
GP = anna .
\end{lstlisting}

\begin{lstlisting}[caption={Вопрос и ответ №5}, label={lst:q5}]
?- find_grandparent(GP, _, female, ivan).
GP = petr ;
GP = anna .
\end{lstlisting}

\begin{lstlisting}[caption={Вопрос и ответ №6}, label={lst:q6}]
?- max2_no_cut(1, 2, Res).
Res = 2.
\end{lstlisting}

\begin{lstlisting}[caption={Вопрос и ответ №7}, label={lst:q7}]
?- max2_cut(3, 1, Res).
Res = 3.
\end{lstlisting}

\begin{lstlisting}[caption={Вопрос и ответ №8}, label={lst:q8}]
?- max3_no_cut(3, 1, 2, Res).
Res = 3.
\end{lstlisting}

\begin{lstlisting}[caption={Вопрос и ответ №9}, label={lst:q9}]
?- max3_cut(0, 1, 2, Res).
Res = 2.
\end{lstlisting}

\begin{lstlisting}[caption={Вопрос и ответ №10}, label={lst:q10}]
?- max3_cut(0, 5, 2, Res).
Res = 5.
\end{lstlisting}
